\section{Further research questions and proposed next investigations}
\label{rq:section}

The following questions begin beyond the proved statements. They seek exact
identities, structural classifications, and finite evaluations. In
particular, a failure to recognize a numerical value in a short constant
basis is not a result about nonreducibility.

\begin{enumerate}[label=\textbf{Q\arabic*.}]
\item \textbf{A short recurrence for the full bivariate contact sequence.}
The coefficient algorithm \eqref{ct:recursion}--\eqref{ct:dk} computes every
contact, and \eqref{ct:topplus}--\eqref{ct:topminus} evaluate their highest
spectral degrees. Find a recurrence in the contact index \(m\), with
polynomial coefficients in \(p,q,a-b\), that preserves every lower-degree
Gamma correction. A satisfactory recurrence should recover the diagonal
Euler-polynomial law and exact cutoff translation without expanding all
intermediate shifted factors.

\item \textbf{Exact factorization and zero loci of the contact polynomials.}
Parity explains some factors of \(c_m,d_m\), but it does not classify all
special parameter values where a contact vanishes. Determine the complete
factorization pattern on \(p+q=-2m\), separating factors forced by exchange
symmetry from additional algebraic factors. The first target is a uniform
description of the residual factors for \(m=4,5,6\), proved from the
coefficient generator rather than inferred from numerical roots.

\item \textbf{An independent outer spectral weight.}
Adjoin \((n+c)^{-w}\) to a product of independently ordered ordinary or
alternating Hurwitz tails. Determine the moving polar hyperplanes and give
a finite-cutoff identity with a precisely specified remainder coordinate.
The present Gamma normalization chooses its outer weight through the
orders \(p,q\); it does not solve the additional independent variable
\(w\). Mellin integration of the completed endpoint law is a natural
starting point, but its boundary terms must be retained.

\item \textbf{Cyclotomic contact matrices.}
For a primitive \(d\)th-root twist with \(d>2\), replace the two-step
alternating recurrence by its \(d\)-component system. Seek a Gamma-product
normalization and an exact matrix-valued bilinear endpoint law. Determine
the resulting resonance lattice, every independent-order contact, and the
effect of shifting the cutoff by one lattice step. The case \(d=3\) is a
concrete first target; the alternating case supplies the normalization test.

\item \textbf{Three genuinely independent tails.}
For three shifts, three orders, and independent colors, classify the
smallest adjacent-product certificate whose difference removes all
unwanted residual terms. If depth-two sums are unavoidable within the
chosen certificate class, give their exact coefficients and endpoint
terms. The aim is a specified finite identity, with the size of its
function alphabet proved within that class; powers of one tail are not
sufficient evidence for the independent-data case.

\item \textbf{One ordinary full-jet integral for odd polynomial weights.}
The compact connected transform already handles arbitrary polynomial
weights, and the full even-weight jets admit the ordinary master derived
here. For an odd weight, produce a comparably economical full-jet formula
that retains all resonant Stieltjes coefficients and the inherited local
correction. A first target is \(W(x)=x\), two masses, and the third jet.
The subtraction must state its independent regulator and agree with
\eqref{sp:inherited-correction}; the known local degree theorem itself is
not the missing result.

\item \textbf{Finite special-value reductions on explicit mass loci.}
At \(a=1/2\), investigate the cubic compact integral
\eqref{da:half-cubic} and the scalar third jet at a mass satisfying
\(e^{-2\pi\sqrt\mu}=1/2\). Seek a proved reduction in a declared alphabet
of classical or multiple polylogarithms, their order derivatives, and
Gamma or zeta derivatives. Keep a finite evaluation separate from the
ordinary integral representation already established. If an obstruction
is proposed, formulate and prove it in a precise function or
difference-equation class.

\item \textbf{Complex masses and controlled branch crossings.}
Continue the density formulas to nonreal masses with explicitly chosen
contours, including conjugate pairs. Determine exact jump formulas when a
mass crosses a negative lattice value \(- (n+a)^2\). The expected jump
should be expressible through the affected finite spectral terms, but
the contour deformation and spectral normalization must be proved
compatible. This would extend the real-mass results to a useful complex
parameter geometry.

\item \textbf{Independent rational-shift projections.}
Apply finite residue-class or character projections to the all-order
digamma kernels, and compare the resulting Dirichlet-\(L\) derivatives
with the polylogarithmic density transform. Produce a complete dictionary
at conductors three and four, including the shift primitives and
coincident-mass coefficients. Specify both parity components; a
reflection identity for their sum does not determine them separately.

\item \textbf{Arithmetic reduction of selected cubic Stieltjes coefficients.}
Starting from the already-proved finite Tornheim germ, compute an explicit
small-index table containing
\(\FP\int_0^1\gamma_1^2\gamma_0'\) and its first independent
argument-derivative companions. Determine all relations obtainable from
shuffle, stuffle, symmetry, and differentiation in a declared formal
alphabet. The unresolved target is a further arithmetic reduction or an
exact formal dimension statement, not existence of the generating germ.

\item \textbf{The full reconstruction and gap-six problems.}
Retain the manuscript's unrestricted unequal-even-order reconstruction
question as a separate structural problem. A useful intermediate result
would be a bordered-Hankel identity for the entire gap-six family,
including the moving \(B_0,B_1\) terms and omitted-binomial conventions.
The known eventual nonvanishing mechanism and finite-order completions
must both be recovered. Another finite table would not resolve the
uniform question.

\item \textbf{Formal proofs and validated evaluation of the new masters.}
Formalize first the finite endpoint telescope, the terminating contact
coefficient algebra, and the polynomial subtraction identities.
Then formalize the uniform asymptotic remainder and the two Cauchy
representations with their absolute-Fubini arguments. Separately build
interval evaluators that split every logarithmic endpoint and use
analytically controlled tails. Exact polynomial checks, floating-point
diagnostics, and analytic proof have different roles and should remain
separate in the repository.
\end{enumerate}

The most direct continuation of the present work is Q6: it asks for the
remaining parity sector of an already explicit full-jet representation.
Q1 and Q4 develop the completed contact algebra structurally. Q7 and Q10
instead aim at finite arithmetic reductions, where the existing exact
integral and coefficient formulas provide reliable starting values and
normalization constraints.
